\documentclass[12pt]{elsarticle}

\usepackage{amsmath}
\usepackage{amssymb}
\usepackage{amsfonts}
\usepackage{graphicx}
\usepackage[T1]{fontenc}
\usepackage[utf8]{inputenc}
\usepackage{booktabs}
\usepackage{geometry}
\usepackage{setspace}
\usepackage{threeparttable}
\usepackage{natbib}
\usepackage[hyphens]{url}
\usepackage{xurl}
\usepackage{hyperref}
\hypersetup{
    colorlinks=true,
    linkcolor=black,
    citecolor=black,
    filecolor=magenta,
    urlcolor=blue,
    pdftitle={The Price of Exclusion: SME Set-Asides in Public Procurement},
    pdfauthor={Darcio Genicolo-Martins}
}

\geometry{left=28mm,right=30mm,top=28mm,bottom=28mm}
\doublespacing
\biboptions{authoryear}
\setlength{\bibsep}{0pt plus 0.3ex}
\setlength{\emergencystretch}{3em}

\graphicspath{{figures/}{../output/figures/}{../../v6-jpube/output/figures/}}

% v8 should eventually have its own generated macro registry. During drafting,
% fall back to the v6 registry so section drafts can reuse existing magnitudes.
\InputIfFileExists{../output/values.tex}{}{%
  \InputIfFileExists{../../v6-jpube/output/values.tex}{}{}}

\begin{document}

\begin{frontmatter}

\title{\textbf{The Price of Exclusion:\\ SME Set-Asides in Public Procurement}}

\author[INSPER]{Darcio Genicolo-Martins}\ead{darciogm1@insper.edu.br}
\address[INSPER]{INSPER, S\~ao Paulo}

\begin{abstract}
SME set-asides support protected firms by excluding their rivals. This paper
shows why that design is costly. I study a legal reinterpretation that expanded
SME-only tendering in S\~ao Paulo's electronic procurement platform and use
reverse-auction drop-out bids to recover type-specific cost distributions. The
set-aside changes two margins with opposite signs: it removes non-SME bidders
from the price-forming pool, but it also induces SME entry and changes SME
composition. A structural decomposition shows that entry responds, but
exclusion dominates. Removing non-SMEs raises simulated prices by
\bneEffectIntensNp/\bneEffectIntensPh{} of the reference price in
non-pharmaceutical/pharmaceutical procurement; post-policy SME entry offsets
\bneEffectEntryNp/\bneEffectEntryPh, leaving positive net effects. Welfare
losses equal \welfLossPctLthirtyNp/\welfLossPctLthirtyPh{} percent of the
open-regime price at $\lambda=\lambdaMain$. A 10 percent SME price preference
preserves non-SME participation and has near-zero price effects in thick
standardized markets, but delivers less redistribution. The policy implication
is a frontier, not a ban on SME support: use exclusion only when the planner is
prepared to pay for the missing competitive discipline.
\end{abstract}

\begin{keyword}
public procurement \sep SME set-asides \sep auctions \sep structural estimation \sep welfare
\MSC[JEL] H32 \sep H57 \sep L26 \sep L53 \sep D44
\end{keyword}

\end{frontmatter}

\section{Introduction}

Set-asides open SME access to public contracts by closing the auction to
everyone else. The design is mechanical: protected firms come in, unprotected
firms go out. Removing rival bidders weakens the competitive discipline that
forms the price. Admitting more SMEs into a now-restricted contest pushes
back. The welfare case depends entirely on which force is larger, and on how
the protected SME pool reorganizes once the auction is closed to non-SMEs.
The empirical literature has documented the cost of SME-favoring rules but
rarely separates these two channels, even though policy design lives in their
balance.

I show that a set-aside replaces competition with eligibility, and that when
the bidders it removes are the price-forming ones the replacement is partial,
expensive, and dominated by a price preference. The point is not that SME
set-asides can be costly: \citet{marion2007} and \citet{nakabayashi2013}
document that they often are. \citet{athey2013} and
\citet{krasnokutskayaseim2011} go further and show that the design of SME
support shapes whether the cost falls on allocation, entry, or both. None of
these decomposes the price effect into what an exclusionary rule does to the
price-forming pool and what the protected pool does in response. Without that
split, the policy debate runs as SME support versus no SME support. The
relevant comparison is different: support that preserves the price-forming
pool against support that achieves access by deleting rival bidders. This
paper places a real procurement reform inside that comparison.

The setting is S\~ao Paulo's centralized electronic procurement platform, the
Bolsa Eletr\^onica de Compras (BEC). In late 2017 and early 2018, a legal
reinterpretation moved the state's SME-only eligibility threshold from the
total value of a purchase notice to the item level. The change mechanically
pushed medical and hospital supplies (Group 65) into SME-only tendering,
because such purchases combine many small items inside large notices. The
empirical cutoff is \policyCutoffMonth{}, when Group-65 public buyers began
mass adoption of the new functionality. The trigger was legal and
administrative, not a medical-supply-specific policy, and it changed
admissibility rather than the production technology of the goods.

The first-stage facts fit the institutional account. After the cutoff,
SME-only adoption in Group 65 rises sharply, non-SME participation falls, SME
participation roughly doubles, and winning prices rise. A compact
difference-in-differences benchmark against \controlGroupCount{} never-treated
product groups places the associated price movement at
\didLowerPct--\didUpperPct{} percent in the structural window. I use that as
a sign, timing, and scale check, not as the main estimate. The central object
is a structural decomposition, under explicit auction assumptions, of how
price formation moves when the policy removes non-SMEs and reshuffles the
active SME pool.

The auction format makes the decomposition feasible. BEC's main procurement
format is the electronic \emph{Preg\~ao}, an English-reverse auction in which
firms submit successively lower offers until one remains. Under independent
private values, losing bidders' drop-out prices reveal their costs
\citep{haile2003}. That is the identifying advantage of the setting and also
its main structural restriction. I use drop-outs to recover type-specific cost
distributions for SMEs and non-SMEs, correct for auction-level scale
heterogeneity in the spirit of \citet{krasnokutskaya2011}, and simulate
counterfactual auctions under observed equilibrium entry. The exercise is not
a full entry equilibrium and not a reduced-form treatment-effect estimate:
observed pre- and post-policy entry rates are taken as primitives, and the
post-policy SME cost distribution is estimated from post-policy exits rather
than derived from a formal selection model.

The decomposition uses three counterfactual prices. Let $S_1$ be the open
auction with the pre-policy bidder pool, $S_2$ the SME-only auction holding
the pre-policy SME pool fixed, and $S_3$ the SME-only auction with the
observed post-policy SME pool. Then
\[
    p_{S_3} - p_{S_1}
    =
    \underbrace{p_{S_2} - p_{S_1}}_{\text{lost competitive discipline}}
    +
    \underbrace{p_{S_3} - p_{S_2}}_{\text{protected-pool offset}}.
\]
The first term measures the cost of removing non-SMEs from the price-forming
pool. The second measures the offset created by the post-policy protected
pool, which combines additional SME participation with changes in the
composition of active SME bidders. The decomposition turns a set-aside price
effect into a design question: how much of the cost comes from replacing
competition by eligibility, and how much is recovered by the protected pool
that the policy brings in?

The main result is that the protected pool responds, but exclusion dominates.
In non-pharmaceutical supplies, removing non-SMEs while holding the pre-policy
SME pool fixed raises simulated prices by \bneEffectIntensNp{} of the
reference price; the observed post-policy SME pool offsets that shock by
\bneEffectEntryNp{}, leaving a net effect of \bneEffectTotalNp{}, with the
exclusion channel accounting for \bneShareIntensNp{} percent of the two
components in absolute magnitude. The protected pool attracts and reshuffles
SMEs, but it does not recreate the competitive discipline supplied by
excluded non-SMEs. Pharmaceutical procurement shows the same qualitative
pattern at larger magnitudes, with composition turnover that makes the
welfare ranking conditional rather than headline.

This price-formation result changes the policy comparison. A full set-aside
supports SMEs by removing their rivals; a price preference does so while
keeping non-SMEs in the auction. As a static design benchmark, I use a 10
percent SME price preference under fixed entry and cost primitives. In thick,
standardized non-pharmaceutical markets, the preference has near-zero price
and welfare cost while the full set-aside generates a substantial loss. The
preference is not a free version of the set-aside; it delivers smaller
redistribution. The comparison is therefore a frontier, not a ranking: a
planner should choose the set-aside only when the additional SME surplus is
worth the allocative and fiscal cost of removing non-SMEs.

The strongest policy result is in thick standardized markets. In thinner and
more heterogeneous pharmaceutical markets, the ranking depends on how the
post-policy SME pool is modeled. The main specification estimates the
post-policy SME cost distribution from post-policy exits; a strict invariance
benchmark instead holds the SME cost distribution fixed across the cutoff.
The non-pharmaceutical ranking is stable across these readings; the
pharmaceutical ranking is not. I treat the gap as a scope condition rather
than a second headline: bidder exclusion is hardest to defend where the
protected pool is thick enough for a preference rule to tilt allocation
without removing competitive discipline, and most defensible, if at all,
where the protected pool is thin, the good is heterogeneous, and the planner
places high value on SME surplus.

The paper develops that frontier through layered analysis built around a
single engine. The engine is the decomposition: the set-aside's price effect
splits into a discipline-loss channel from removing non-SMEs and a
protected-pool offset from the SME response. Drop-out bids in the electronic
Preg\~ao give the channels empirical content without invoking bid-function
inversion \citep{krasnokutskaya2011,krasnokutskayaseim2011}, and a planner
problem then prices the channels against transfers to SMEs, allocative waste,
and the marginal cost of public funds \citep{saezstantcheva2016}. The lesson
is direct: when the bidders a set-aside removes are the ones forming the
price, the rule trades competitive discipline for an incomplete protected-pool
response and earns the difference as a welfare loss. Full set-asides remain
defensible, but the defense requires the planner to assign explicit weight to
SME surplus, market access, or dynamic capacity gains beyond what the
protected pool itself delivers. The natural benchmark for SME procurement
support is not no support, but support that keeps rival bidders inside the
auction.

\section{Setting, Data, and First-Stage Facts}\label{sec:setting}

The empirical episode is a legal change that shifted which firms could compete
for eligible Group~65 items. Around the same cutoff, SME-only take-up rises,
non-SME participation falls, SME participation rises, and prices move up. These
facts establish timing and direction. They do not, by themselves, identify the
structural price mechanism.

\subsection{Legal shock}

Brazil's federal SME statute, \LCSmeStatute{}, requires exclusive SME tenders
for procurement items valued at R\$\,\setAsideThreshBRL{} or less. The 2014
amendment, \LCSetAsideAct{}, made the SME-only rule mandatory rather than
optional. For most goods on S\~ao Paulo's BEC platform, that rule already
operated before the period studied here.

Group~65---medical, dental, and hospital equipment and supplies---was different
for a legal reason. The prevailing interpretation applied the
R\$\,\setAsideThreshBRL{} threshold to the total value of the purchase notice,
not to each item or lot. Because medical-supply notices often bundle many small
items inside large purchase orders, the total-notice interpretation kept many
individually eligible items outside the SME-only default.

The relevant change was a generic legal reinterpretation. BEC enabled
SME-only purchase-order functionality in COMUNICADO BEC
n\textordmasculine~\comunicadoBECone{} on 18 July 2017 and mixed-exclusivity
notices in COMUNICADO BEC n\textordmasculine~\comunicadoBECtwo{} on 20
November 2017. PGE-SP then issued Parecer Sub-G Cons.\
n\textordmasculine~\pareceNum{} on 11 December 2017, holding that the
R\$\,\setAsideThreshBRL{} threshold applies item-by-item regardless of the
total value of the purchase offer. The empirical cutoff is
\policyCutoffMonth{}, when Group~65 public buyer units began mass take-up of
the functionality. TCE-SP aligned with the new interpretation on 16 May 2018,
after the empirical cutoff, ratifying rather than initiating the operational
sequence.

The institutional claim is narrow. The trigger was not a Group~65 market
intervention and did not speak to medical-supply prices, supplier complaints,
or technology. It was a general reading of procurement law that interacted with
the way Group~65 notices were written. The cutoff shifts admissibility and the
active bidder pool; it does not by itself identify a full entry model or rule
out composition changes among SMEs.
Online Appendix~OA-A reports the institutional timeline and platform sample
details.

\subsection{Data and auction format}

The data are administrative microdata from BEC, S\~ao Paulo's centralized
electronic procurement platform. The raw extract contains \nBidsRaw{} million
bid-level observations across \nItemsRaw{} items, \nPurchaseOrders{} purchase
orders, and \nPbu{} public buyer units. Group~65 accounts for roughly
\groupSixtyfivePctTx{} percent of platform transactions in the window used here
and is the treated group. The reduced-form benchmark compares Group~65 to
\controlGroupCount{} never-treated product groups.

The structural exercise uses Group~65 Preg\~ao auctions in a symmetric
\structuralWindowM-month window on each side of the \policyCutoffMonth{} cutoff,
from \windowStart{} through \windowEnd{}. The Preg\~ao format is central: it is
an electronic English-reverse auction in which the auction opens at a buyer
reference price, firms submit lower offers, and firms exit as the clock falls.
Under the independent-private-values interpretation used in the model, losing
drop-out prices are cost observations.

For that reason, the structural sample restricts attention to Preg\~ao auctions
with at least \filterMinBidders{} active firms and normalized bids
$c_\epsilon=b/p^{\mathrm{ref}}\in(0,\filterCepsUpper]$. The resulting sample
contains \nObsAuctions{} firm-auction observations across
\nDistinctAuctions{} auctions, split almost evenly between the pre-period
(\nPreAuctions{} auctions) and post-period (\nPostAuctions{} auctions).

Each observation is classified by bidder type and product class. SME status is
taken from the BEC administrative registry and validated against historical
firm bidding records. I split Group~65 into non-pharmaceutical supplies and a
pharmaceutical class covering CMED-regulated medicines and biologicals. This
split is not a second treatment. It is a market-thickness distinction: the
non-pharmaceutical segment is closer to a thick standardized market, while
pharmaceutical procurement has fewer SME bidders and more within-class
heterogeneity. The main policy claim is strongest in the former; the latter is
used as a boundary case.

\subsection{First-stage facts}

Four facts matter. SME-only adoption rises in Group~65 after the cutoff;
state-level adherence reaches \adherenceGsixtyfivePost{} percent rather than
full coverage. Non-SME participation falls. SME participation rises. Winning
prices move up in the same period. The participation responses are the
reduced-form moments most useful for identification because they do not pass
through the reference-price normalization; the price movement that follows
is taken as directional validation, with magnitude estimated structurally in
Section~\ref{sec:price_formation}.

Table~\ref{tab:first_stage_v8} reports the participation margins in the
structural sample. In non-pharmaceutical Preg\~ao auctions, the average number
of SME bidders rises from \bneNsmePreNp{} to \bneNsmePostNp{}, while non-SME
participation falls from \bneNnsPreNp{} to \nonSmePostNp{}. In pharmaceutical
auctions, SME participation rises from \bneNsmePrePh{} to \bneNsmePostPh{},
while non-SME participation falls from \bneNnsPrePh{} to \nonSmePostPh{}. The
post-period non-SME counts remain positive because take-up is incomplete, but
both margins move in the direction required by the decomposition.

\begin{table}[!htbp]
\centering
\begin{threeparttable}
\caption{First-stage participation facts in Group~65 Preg\~ao auctions}
\label{tab:first_stage_v8}
\small
\setlength{\tabcolsep}{6pt}
\begin{tabular}{lcc}
\toprule
 & Non-pharma & Pharma \\
\midrule
Pre: SME bidders per auction & \bneNsmePreNp{} & \bneNsmePrePh{} \\
Pre: non-SME bidders per auction & \bneNnsPreNp{} & \bneNnsPrePh{} \\
Post: SME bidders per auction & \bneNsmePostNp{} & \bneNsmePostPh{} \\
Post: non-SME bidders per auction & \nonSmePostNp{} & \nonSmePostPh{} \\
\addlinespace
Change in SME bidders & +0.93 & +0.67 \\
Change in non-SME bidders & -1.18 & -0.95 \\
\bottomrule
\end{tabular}
\begin{tablenotes}\footnotesize
\item Sample: Group~65 Preg\~ao auctions in the \structuralWindowM-month window
on each side of the \policyCutoffMonth{} cutoff, after the structural filters
$c_\epsilon\in(0,\filterCepsUpper]$ and at least \filterMinBidders{} active
firms. Counts are average distinct bidders per auction by type and class.
Post-period non-SME counts remain positive because the empirical SME-only
take-up rate is \adherenceGsixtyfivePost{} percent, not 100 percent.
\end{tablenotes}
\end{threeparttable}
\end{table}

The price-forming order statistic moves in the same direction. The DiD
specification applied to $\log b^{(2)}$, the second-lowest bid in the
iterative Preg\~ao phase and the price-forming object in the IPV
interpretation, delivers \bidSecondDidEstNp{} (SE \bidSecondDidSeNp{}) in
non-pharmaceuticals and \bidSecondDidEstPh{} (SE \bidSecondDidSePh{}) in
pharmaceuticals. These estimates are built from raw bid units and do not
pass through the reference-price normalization, so they corroborate the
order-statistic channel directly; Online Appendix~OA-D.4 reports the
full event study.

Winning prices also move. Table~\ref{tab:did_benchmark_v8} reports the DiD
on $\log p^{\mathrm{final}}$ comparing Group~65 to never-treated product
groups, with item and month fixed effects, auction-format and quantity
controls, PBU controls, and item-clustered standard errors. The coefficient
is written in a ``DiD-in-reverse'' convention: a negative coefficient on
Group~65 interacted with the pre-period means that the open regime was
cheaper, so the post-cutoff SME-only extension undoes that discount. Across
windows, the log-price coefficient is between -0.108 and -0.148 with PBU
controls; the \structuralWindowM-month benchmark is \didPriceEighteenPbu{}.
This coefficient is used as a directional and timing check rather than as
the source of the structural magnitude; the latter is built in
Section~\ref{sec:price_formation} in $p^{\mathrm{ref}}$-normalized units, and
Section~\ref{sec:pref_endog} examines reference-price endogeneity directly.

\begin{table}[!htbp]
\centering
\begin{threeparttable}
\caption{Reduced-form price benchmark}
\label{tab:did_benchmark_v8}
\small
\setlength{\tabcolsep}{7pt}
\begin{tabular}{lccc}
\toprule
Window around cutoff & 6 months & 12 months & 18 months \\
\midrule
$g65 \times Pre$ on $\log p^{\mathrm{final}}$ & -0.148 & -0.108 & \didPriceEighteenPbu{} \\
Standard error & (0.014) & (0.013) & (0.012) \\
Observations & 219,535 & 439,054 & 649,714 \\
\bottomrule
\end{tabular}
\begin{tablenotes}\footnotesize
\item Estimates use item and month fixed effects, controls for sealed-bid
format and log quantity, PBU controls, and item-clustered standard errors.
The coefficient is reported in the same sign convention: a negative
pre-period coefficient is the open-regime price discount that disappears after
the SME-only extension. The DiD is used as a sign, timing, and scale benchmark;
the structural sections estimate the mechanism.
\end{tablenotes}
\end{threeparttable}
\end{table}

The reduced form has a limited role: the same legal shock changes
admissibility, the active bidder pool, and prices in a single step.
Section~\ref{sec:identification_architecture} sets out how the structural
exercise separates those margins; OA-B reports the balance, placebo, and
alternative-DiD diagnostics behind the conservative reading.

\subsection{Identification architecture}\label{sec:identification_architecture}

The argument rests on three pieces of evidence that carry separate identifying
weight. The difference-in-differences against \controlGroupCount{}
never-treated groups validates a real first-stage footprint at
\policyCutoffMonth{}; the structural decomposition, conditional on observed
entry and recovered costs, identifies the price-forming mechanism; the static
welfare comparison of $V_0$ against $V_3$, conditional on recovered
primitives, identifies the policy implication. None of the three inherits
fragilities from the others. Each piece can fail in a different way without
compromising the rest.

The reduced form requires only that the legal episode be the dominant shock
to Group~65 around \policyCutoffMonth{}. Modern estimators reduce magnitude
relative to TWFE (\didAltCsEst{} versus \didAltTwfeEst{}), which is consistent
with the role assigned: timing and sign, not structural magnitude. The
structural decomposition does not require cross-group balance, because it
conditions on within-period type-specific cost distributions and equilibrium
bidder counts; it requires the IPV interpretation of Preg\~ao drop-outs,
disciplined in Section~\ref{sec:robustness}. The welfare comparison requires
neither balance nor the DiD; it requires that the recovered primitives be
interpretable, which is what the same section examines.

Five concerns recur. Pre-treatment imbalance between Group~65 and the pooled
controls limits the force of the reduced form (OA-B reports the balance
table); it does not affect the structural exercise, which uses no
cross-group comparison. Modern DiD attenuation is consistent with the
conservative reading the paper assigns to the reduced form (\didAltCsEst{}
versus \didAltTwfeEst{}, OA-B). Reference-price endogeneity is examined directly in
Section~\ref{sec:pref_endog} and Online Appendix~OA-D.4: a static DiD on
$\log p^{\mathrm{ref}}$ moves by \pRefDidEstNp{} (NP) and \pRefDidEstPh{}
(PH), well below the corresponding DiD on $\log p^{\mathrm{final}}$, and an
internal placebo confirms the headline price effect survives where
$p^{\mathrm{ref}}$ is stable.
IPV is the maintained restriction, disciplined in
Section~\ref{sec:robustness} and OA-C through censoring, common-scale, and
cross-modality diagnostics. The empirical cutoff is the mass take-up date
rather than a legal moment; Section~\ref{sec:setting} documents the
institutional sequence, and the phased-adoption sensitivity in OA-B excludes
platform enablement months and retains the qualitative result.

The exercise does not identify a free-entry primitive, a structural model of
SME selection, or dynamic capacity gains. Entry rates are observed
equilibrium objects in the spirit of \citet{atheyseira2011}, and the
post-policy SME cost distribution is an observed active-bidder distribution
with strict invariance reported as a bounding exercise in
Section~\ref{sec:robustness}.

\section{Model and Identification}\label{sec:model}

The empirical facts in Section~\ref{sec:setting} show that the reform moves
participation and prices in the direction implied by the institutional shock.
This section defines the structural objects used to separate the two channels.
The model is intentionally spare. It uses the auction format to recover
type-specific costs, takes observed entry as an equilibrium object, and
simulates three counterfactual bidder pools. The goal is not to estimate a full
entry equilibrium; it is to decompose the price effect of exclusion into lost
competitive discipline and the offset created by protected entry and
composition.

\subsection{Auction environment}

Consider a single procurement auction for one indivisible contract. Each firm
$i$ has type $k_i\in\{\mathrm{SME},\neg\}$ and draws a private supply cost
$c_i$ from a type-specific distribution $F_c^k$. The two distributions are
allowed to differ. The buyer runs an electronic Preg\~ao, an English-reverse
auction in which the price clock moves downward from the buyer's reference
price and firms remain active until they are no longer willing to supply at the
current price. The last active firm wins and supplies the contract at the
clearing price.

Under the independent-private-values interpretation, exit at cost is weakly
dominant in this format \citep{haile2003}. A losing firm's drop-out price is
therefore a direct observation of its cost. This is the key identifying
advantage of the setting. In first-price procurement auctions, the researcher
observes strategically shaded bids and must recover costs by inverting a bid
function. Here the price-forming statistic can be recovered from exits. The
winner is censored because the winner does not drop out; its cost is known only
to be no higher than the second-lowest exit. The baseline estimator includes
the winner's final price as a tight upper-bound observation, while the
robustness section reports a Turnbull NPMLE that treats the winner as
left-censored at that bound.

The expected transaction price in the clock interpretation is the expected
second-lowest cost, $E[c_{(2)}]$, while the production cost of the allocation is
the winner's cost, $E[c_{(1)}]$. This distinction matters later for welfare:
$c_{(2)}$ determines the government's payment, while $c_{(1)}$ determines
allocative efficiency.

\subsection{Recovering cost primitives}

The structural sample is organized by class, period, and bidder type. For each
cell, loser drop-outs identify the corresponding empirical distribution of
normalized costs. Costs are normalized by the buyer's reference price, so the
object entering the simulation is a distribution of $c_i/p^{\mathrm{ref}}$.
Reference-price normalization absorbs the main cross-item scale differences,
but it does not remove all auction-level heterogeneity. Two auctions in the
same product class can differ in volume, delivery burden, regulatory
documentation, or item-specific complexity in ways that move every bid in the
auction together.

I therefore use the standard multiplicative auction-level heterogeneity
correction of \citet{krasnokutskaya2011}. Let
\[
    y_{it}=\log(b_{it}/p_t^{\mathrm{ref}})=a_t+e_{it},
\]
where $a_t$ is an auction-level scale component and $e_{it}$ is the
bidder-specific cost component. A method-of-moments variance decomposition
estimates the between-auction and within-auction components, and a
best-linear-predictor shrinkage removes $\hat a_t$ from each observation. The
simulation uses the cleaned residual cost distribution
$\exp(\hat e_{it})$. In the v6 estimates, the resulting Preg\~ao intraclass
correlations range from \iccNpSme{} to \iccPharmaNonSme{} in the main cells,
with values comparable to procurement settings where this correction is
standard.

Two diagnostics discipline the cost-recovery step. First, in a pre-policy cell
unaffected by the set-aside, the cost distribution recovered from Preg\~ao
drop-outs can be compared to a GPV recovery from Convite first-price bids
\citep{gpv2000}. The two recoveries are close in the pharma non-SME pre-period
cell, which is the cleanest falsification cell because those firms are not
targeted by the policy. Second, the Turnbull NPMLE for winner censoring delivers
the same decomposition qualitatively: the within-auction share remains large
(\turnbullShareNp{} percent in non-pharma and \turnbullSharePh{} percent in
pharma), and the total effect differs from the baseline by at most
\turnbullDeviationPct{} percent across the reported cells. These checks do not
prove IPV. They show that the main cost-recovery step is not being driven by
winner censoring or by an obvious conflict between auction formats.

\subsection{Entry and the post-policy SME distribution}

Entry is treated as observed equilibrium entry. For each class and period, I
use the observed average numbers of SME and non-SME bidders per auction as the
arrival rates for the simulation. In the open pre-policy regime, both SMEs and
non-SMEs may enter. In the SME-only counterfactual, non-SMEs are removed from
the eligible pool. The observed post-policy SME count is used for the
post-policy SME-only scenario.

This choice is deliberate. A full entry model would require assumptions on
fixed costs, information, bidder-specific participation decisions, and how the
legal shock changes expected profits. The legal shock directly identifies
admissibility and observed participation, not a free-entry primitive. The paper
therefore follows the reduced-form-entry spirit of
\citet{atheyseira2011}: observed pre- and post-policy bidder counts are
equilibrium objects used in counterfactual price formation.

The same caution applies to the post-policy SME cost distribution. The baseline
simulation estimates $F_c^{\mathrm{SME,Post}}$ directly from post-policy
drop-outs and uses it as the cost distribution of active post-policy SME
bidders. This object may reflect selection into the protected market,
technology or sourcing changes within the SME segment, or other composition
changes around the cutoff. The baseline does not derive that difference from a
formal selection model. The strict-invariance benchmark instead imposes
$F_c^{\mathrm{SME,Post}}=F_c^{\mathrm{SME,Pre}}$. The two readings bracket how
much the policy conclusion depends on the empirical post-policy SME
distribution. The non-pharmaceutical policy ranking is stable across them; the
pharmaceutical ranking is not, which is why pharma is treated as a boundary
case rather than as the headline.

\subsection{Counterfactual objects}

The decomposition uses three simulated auction environments:
\begin{itemize}
    \item $S_1$: open auction, pre-policy bidder pool. SMEs and non-SMEs enter
    at their observed pre-policy rates, with pre-policy cost distributions.
    \item $S_2$: SME-only auction, pre-policy SME pool. Non-SMEs are removed,
    while the SME arrival rate and SME cost distribution are held at their
    pre-policy values.
    \item $S_3$: SME-only auction, post-policy SME pool. Non-SMEs are removed,
    and SMEs enter with the observed post-policy arrival rate and the baseline
    post-policy SME cost distribution.
\end{itemize}

For each scenario, I draw bidder counts from the corresponding type-specific
arrival process, draw costs from the recovered cost distributions, sort the
draws, and record $c_{(1)}$ and $c_{(2)}$. The simulated price is
$p_S=E_S[c_{(2)}]$, normalized by the reference price. The simulated production
cost is $E_S[c_{(1)}]$.

The total set-aside price effect is
\[
    \Delta^{SA}=p_{S_3}-p_{S_1}.
\]
It decomposes mechanically into two terms:
\begin{equation}\label{eq:main_decomp}
    p_{S_3}-p_{S_1}
    =
    \underbrace{p_{S_2}-p_{S_1}}_{\Delta^{excl}
    \text{ lost competitive discipline}}
    +
    \underbrace{p_{S_3}-p_{S_2}}_{\Delta^{entry}
    \text{ entry/composition offset}}.
\end{equation}
The first term holds the SME pool fixed and removes non-SMEs. It captures the
price-forming cost of bidder exclusion: the second-order statistic is now drawn
from the protected pool rather than from the full pool. The second term allows
the post-policy SME pool to enter. It captures the extent to which protected
entry and composition offset the exclusion shock. This term can be negative,
and in the data it is: additional SME participation partially restores
competition, but not enough to undo the price increase.

The labels in equation~\eqref{eq:main_decomp} are intentionally descriptive
rather than behavioral. In an English-reverse IPV auction, surviving bidders do
not need to shade bids differently when rivals are excluded; exit at cost
remains weakly dominant. The lost-discipline component is therefore an
order-statistic effect of admissibility, not evidence that SMEs bid less
aggressively. This is the distinction between the mechanism in this paper and
the first-price settings in which bidder preferences combine cost differences,
order statistics, and strategic bid shading.

\subsection{Identification threats and what disciplines them}

Four threats matter for the interpretation. The first is the IPV/drop-out
restriction. If drop-outs are not informative about costs, the recovered
$F_c^k$ distributions are not structural cost primitives. The cross-modality
check and the Turnbull censoring exercise are the main diagnostics, and Section
6 summarizes an affiliation sensitivity that allows within-auction cost
correlation up to $\rho_c=\ipvRhoMax$.

The second threat is primitive stability. The institutional shock shifts
admissibility, but the recovered SME distribution changes across periods. The
baseline treats that empirical post-policy distribution as the relevant active
SME pool; strict invariance imposes no change. The comparison between those
readings is central to the scope conditions.

The third threat is entry aggregation. The baseline uses average
class-by-period arrival rates. That is enough for the main accounting exercise,
but it is coarse. Section 6 therefore treats empirical entry-count
distributions and finer arrival specifications as robustness checks rather than
as new mechanisms.

The fourth threat is coordination. Bidder-pair screens detect clustering in
the data, so the interpretation of drop-out prices as costs requires care. The
important diagnostic is whether clustering intensifies after non-SMEs are
removed, and whether the decomposition survives in low-coordination samples.
The paper treats this as a first-order threat to interpretation, not as a
cosmetic robustness exercise.

\section{Exclusion, Entry, and Price Formation}\label{sec:price_formation}

The set-aside changes two margins at once. It removes non-SMEs from eligible
auctions, but it also attracts more SMEs into the protected pool. This section
uses the counterfactuals from Section~\ref{sec:model} to separate the two
forces. The result is simple: SME entry responds strongly, but the price effect
is dominated by the loss of non-SME competitive discipline.

\subsection{Entry responds}

The first-stage participation facts already rule out a simple story in which
prices rise because protected firms fail to enter. In non-pharmaceutical
Preg\~ao auctions, the average number of SME bidders rises from
\bneNsmePreNp{} to \bneNsmePostNp{} per auction after the cutoff. In
pharmaceutical auctions, it rises from \bneNsmePrePh{} to \bneNsmePostPh{}.
These are large changes relative to the pre-policy protected pools. Non-SME
participation moves in the opposite direction, falling from \bneNnsPreNp{} to
\nonSmePostNp{} in non-pharma and from \bneNnsPrePh{} to \nonSmePostPh{} in
pharma.

That participation response matters for interpretation. If one held the SME
pool fixed at its pre-policy size, the set-aside would look more expensive than
the realized post-policy environment. The government gets some competition back
through induced SME participation. The key question is whether that offset is
large enough to replace the price discipline supplied by the excluded non-SMEs.

\subsection{Exclusion dominates}

Table~\ref{tab:price_decomp_v8} reports the three simulated price statistics.
The open pre-policy auction is $S_1$. The SME-only auction holding the
pre-policy SME pool fixed is $S_2$. The SME-only auction with the observed
post-policy SME pool is $S_3$. Prices are normalized by the buyer's reference
price and correspond to the simulated second-order statistic, $E[c_{(2)}]$.

The first move, $S_1\rightarrow S_2$, isolates bidder exclusion. In
non-pharmaceutical markets, removing non-SMEs while holding the pre-policy SME
pool fixed raises the simulated price from \bneMeanSoneNp{} to
\bneMeanStwoNp{}, an increase of \bneEffectIntensNp{} of the reference price.
In pharmaceutical markets, the same move raises the simulated price from
\bneMeanSonePh{} to \bneMeanStwoPh{}, an increase of \bneEffectIntensPh{}.
This is the price-forming cost of replacing the full bidder pool with the
protected pool.

The second move, $S_2\rightarrow S_3$, allows the post-policy SME pool to
enter. This term is negative in both classes. In non-pharma, SME
entry/composition reduces the simulated price by \bneEffectEntryNp{} of the
reference price. In pharma, it reduces the simulated price by
\bneEffectEntryPh{}. The realized set-aside effect, $S_3-S_1$, remains
positive: \bneEffectTotalNp{} in non-pharma and \bneEffectTotalPh{} in pharma.
Entry attenuates the price shock; it does not eliminate it.

\begin{table}[!htbp]
\centering
\begin{threeparttable}
\caption{Price decomposition: exclusion versus entry/composition}
\label{tab:price_decomp_v8}
\footnotesize
\setlength{\tabcolsep}{4pt}
\begin{tabular}{lccccccc}
\toprule
 & \multicolumn{3}{c}{Simulated price} & \multicolumn{4}{c}{Decomposition} \\
\cmidrule(lr){2-4}\cmidrule(lr){5-8}
Class & $S_1$ & $S_2$ & $S_3$ &
$S_2-S_1$ & $S_3-S_2$ & $S_3-S_1$ & Excl./net \\
\midrule
Non-pharma
& \bneMeanSoneNp{} & \bneMeanStwoNp{} & \bneMeanSthreeNp{}
& +\bneEffectIntensNp{} & \bneEffectEntryNp{}
& +\bneEffectTotalNp{} & \bneIntensRelNetNp{}\% \\
Pharma
& \bneMeanSonePh{} & \bneMeanStwoPh{} & \bneMeanSthreePh{}
& +\bneEffectIntensPh{} & \bneEffectEntryPh{}
& +\bneEffectTotalPh{} & \bneIntensRelNetPh{}\% \\
\bottomrule
\end{tabular}
\begin{tablenotes}\footnotesize
\item $S_1$ is the open auction with the pre-policy bidder pool. $S_2$ is the
SME-only auction holding the pre-policy SME pool fixed. $S_3$ is the SME-only
auction using the observed post-policy SME pool. Prices are simulated
$E[c_{(2)}]$ values normalized by the buyer reference price. ``Excl./net''
reports $(S_2-S_1)/(S_3-S_1)$. The entry/composition term is negative because
the post-policy protected pool partially offsets the exclusion shock.
\end{tablenotes}
\end{threeparttable}
\end{table}

The relative-to-net column is intentionally reported because it shows why the
absolute share can understate the size of the exclusion margin. In non-pharma,
the exclusion component is \bneIntensRelNetNp{} percent of the net effect; in
pharma it is \bneIntensRelNetPh{} percent. The entry/composition term offsets
\bneEntryRelNetNp{} percent and \bneEntryRelNetPh{} percent of the net effect,
respectively. If instead one normalizes by the sum of absolute component
magnitudes, exclusion accounts for \bneShareIntensNp{} percent in non-pharma
and \bneShareIntensPh{} percent in pharma. Both normalizations point in the
same direction: the dominant force is the loss of the non-SME price-forming
pool, not a lack of SME entry.

Figure~\ref{fig:decomposition_v8} makes the same point visually. The jump from
$S_1$ to $S_2$ is the cost of exclusion. The decline from $S_2$ to $S_3$ is the
offset from the realized protected pool. The remaining gap between $S_1$ and
$S_3$ is the realized price cost of the set-aside under observed entry.

\begin{figure}[!htbp]
\centering
\includegraphics[width=0.95\textwidth]{fig_v3_decomposition.pdf}
\caption[Counterfactual price decomposition.]{\textit{Counterfactual price
decomposition.} Bars report the simulated second-order statistic normalized by
the buyer reference price. $S_1$ is the open pre-policy auction, $S_2$ removes
non-SMEs while holding the pre-policy SME pool fixed, and $S_3$ allows the
observed post-policy SME pool to enter. The movement from $S_1$ to $S_2$ is the
lost competitive-discipline component; the movement from $S_2$ to $S_3$ is the
entry/composition offset.}
\label{fig:decomposition_v8}
\end{figure}

\subsection{Entry and composition are offsets, not the source}

The decomposition separates two facts that are often conflated. The set-aside
does induce protected entry. That entry is valuable in price terms because it
pulls the simulated price down from $S_2$ to $S_3$. But the policy first creates
a larger price increase by removing non-SMEs from the auction. In the baseline
simulation, holding the protected pool at its pre-policy size would make the
latent price shock \latentShockRatioNp{} percent larger in non-pharma and
\latentShockRatioPh{} percent larger in pharma. The realized post-policy SME
pool absorbs part of that shock, but the net effect remains positive in both
classes.

The entry/composition term should not be overinterpreted as pure entry. It also
contains whatever empirical difference exists between the active SME bidders
observed before and after the cutoff. That is why the paper treats
$F_c^{\mathrm{SME,Post}}$ explicitly as an observed post-policy active-bidder
distribution rather than as the output of a formal selection model. Two
diagnostics keep the conclusion focused. First, the decomposition remains
exclusion-dominant under the Turnbull cost estimator, with exclusion shares of
\turnbullShareNp{} percent in non-pharma and \turnbullSharePh{} percent in
pharma on the absolute-share normalization. Second, under the strict-invariance
benchmark that forces the post-policy SME distribution to equal the pre-policy
one, the exclusion component remains the larger component
(\strictShareNp{} percent in non-pharma and \strictSharePh{} percent in
pharma). What changes under strict invariance is the policy ranking in pharma,
not the basic price-formation fact that bidder exclusion is the larger force.

Firm turnover helps explain why the pharma ranking is more fragile. In
pharmaceuticals, \turnoverPhPctNewFirms{} percent of post-policy SME firms are
new relative to the pre-period pool and they account for
\turnoverPhPctNewBids{} percent of post-policy SME bids. In non-pharma, the
corresponding bid share is \turnoverNpPctNewBids{} percent. The composition
shift is therefore larger precisely in the class where the welfare ranking is
model-sensitive. This is a scope condition, not a competing headline. The
robust result is that in the thicker non-pharmaceutical market, the price cost
is driven by removing non-SMEs and survives alternative treatments of the
protected pool.

\subsection{Interpretation}

The mechanism is not that protected firms strategically bid less aggressively
after the set-aside. In the maintained Preg\~ao/IPV model, bidders exit at
cost. The mechanism is that the set-aside changes which cost order statistic
forms the price. In the open auction, non-SMEs remain in the pool and discipline
the second-lowest cost. In the SME-only auction, that discipline is removed.
Additional SME entry pushes back against the loss, but the post-policy
protected pool does not recreate the full-pool benchmark.

This is the price-formation result the welfare section builds on. A policy
that supports SMEs by excluding non-SMEs buys redistribution by weakening the
auction's price-forming pool. A policy that supports SMEs while leaving
non-SMEs in the auction works through a different channel: it can tilt
allocation toward SMEs while preserving the bidders that discipline the price.

\section{Welfare and Policy Design}\label{sec:welfare_policy}

The price decomposition has a direct welfare interpretation. A set-aside raises
government outlays because it changes the order statistic that determines
payment, but the social cost is not the price effect itself. Part of the extra
payment is a transfer to the protected winner; part reflects a real production
inefficiency from assigning the contract to a higher-cost firm; and part is the
tax distortion created by financing the additional public outlay. The policy
question is therefore sharper than whether SME support is desirable. It is
whether the incremental SME surplus generated by excluding rival bidders is
valuable enough to justify the allocative and fiscal cost of removing them from
the auction.

\subsection{Welfare accounting}

Let $p^V$ denote the expected procurement price under policy $V$, and let
$c_{(1)}^V$ denote the winner's production cost. Relative to the open benchmark
$S_1$, the per-auction welfare loss from the SME-only set-aside $V_0$ is
\begin{equation}
    \mathrm{Loss}(V_0)
    =
    \underbrace{c_{(1)}^{V_0} - c_{(1)}^{S_1}}_{
        \mathrm{DWL}_{\mathrm{alloc}}}
    +
    \underbrace{\lambda\left(p^{V_0} - p^{S_1}\right)}_{
        \mathrm{MCPF\ distortion}} .
    \label{eq:welfare_v8}
\end{equation}
The first term is allocative deadweight loss: the additional production cost
created when the restricted auction selects an SME whose cost exceeds the
lowest-cost firm in the full pool. The second term prices the marginal cost of
public funds; the baseline uses $\lambda=\lambdaMain$, consistent with the
standard public-finance range used to evaluate tax-financed spending
\citep{ballard1985,hendren2020,finkelstein2020}. The remaining component of
the government's extra payment is not destruction in this accounting. It is an
implicit transfer to SME bidders, and it becomes a social benefit only to the
extent that the planner places extra welfare weight on SME surplus.

\subsection{The welfare cost of exclusion}

Table~\ref{tab:welfare_policy_v8} summarizes the welfare accounting and the
non-exclusionary policy benchmark. In non-pharmaceuticals, the full set-aside
raises government payments by \welfDeltaGovNp{} of the reference price. Of this,
\welfDwlAllocNp{} is real allocative waste and \welfMcpfDistLthirtyNp{} is the
MCPF distortion at $\lambda=\lambdaMain$, yielding a welfare loss equal to
\welfLossPctLthirtyNp~percent of the open-regime price. In pharmaceuticals, the
same calculation is larger: the payment increase is \welfDeltaGovPh, the
allocative loss is \welfDwlAllocPh, the MCPF distortion is
\welfMcpfDistLthirtyPh, and the total welfare loss is
\welfLossPctLthirtyPh~percent of the open-regime price. The lower endpoints of
the bootstrap intervals remain economically large
(\welfLossPctLoNp~percent in non-pharmaceuticals and \welfLossPctLoPh~percent
in pharmaceuticals), so the welfare conclusion is not a small-sample artifact.

\begin{table}[!htbp]
\centering
\begin{threeparttable}
\caption{Welfare cost of exclusion and the price-preference benchmark}
\label{tab:welfare_policy_v8}
\small
\setlength{\tabcolsep}{5pt}
\begin{tabular}{lcccccc}
\toprule
& \multicolumn{4}{c}{Full SME set-aside $V_0$}
& \multicolumn{2}{c}{10\% preference $V_3$} \\
\cmidrule(lr){2-5}\cmidrule(lr){6-7}
Class
& $\Delta_{\mathrm{gov}}$
& DWL$_{\mathrm{alloc}}$
& MCPF
& Loss / $p_{S_1}$
& $\Delta p$
& SME win gain \\
\midrule
Non-pharma
& \welfDeltaGovNp
& \welfDwlAllocNp
& \welfMcpfDistLthirtyNp
& \welfLossPctLthirtyNp\%
& \optPrefShiftNp
& \prefSmeWinGainTenNp~pp \\
Pharma
& \welfDeltaGovPh
& \welfDwlAllocPh
& \welfMcpfDistLthirtyPh
& \welfLossPctLthirtyPh\%
& \vThreeShiftPh
& \prefSmeWinGainTenPh~pp \\
\bottomrule
\end{tabular}
\begin{tablenotes}\footnotesize
\item $V_0$ is the full SME-only set-aside under the observed post-policy SME
pool. $\Delta_{\mathrm{gov}}$, DWL$_{\mathrm{alloc}}$, MCPF, and $\Delta p$ are
reported in reference-price units. MCPF uses $\lambda=\lambdaMain$. $V_3$ is an
open auction with a 10 percent SME price preference for winner selection; all
firms remain eligible, and the government pays the actual winning bid. SME win
gain is relative to the open benchmark $S_1$.
\end{tablenotes}
\end{threeparttable}
\end{table}

The magnitudes are large enough to matter outside the simulated auction. At the
observed \adherenceGsixtyfivePost~percent Group-65 adherence rate, the
per-auction losses imply an annual welfare cost of roughly
R\$\,\welfAnnualBRLLo~million in S\~ao Paulo Group-65 procurement alone; across
the realistic adherence range of \welfAdherenceLoPct--\welfAdherenceHiPct
percent, the range is R\$\,\welfRangeLoBRL--\welfRangeHiBRL~million. These
figures are not the central identification result, but they discipline the
policy relevance of the mechanism: exclusion is not a symbolic preference. It
is a costly way to transfer surplus.

\subsection{The policy frontier}

The clean benchmark is a price preference rather than no SME support. Under
$V_3$, SMEs receive a 10 percent scoring advantage for winner selection, but
non-SMEs continue to participate and continue to discipline the price-forming
order statistic. The simulated price effect is essentially zero:
\optPrefShiftNp{} of the reference price in non-pharmaceuticals and
\vThreeShiftPh{} in pharmaceuticals. The instrument does not replicate the full
redistribution of $V_0$; the SME win-rate rises by
\prefSmeWinGainTenNp/\prefSmeWinGainTenPh{} percentage points rather than being
forced to 100 percent. That difference is exactly the trade-off the set-aside
hides. Price preferences buy some reallocation toward SMEs while preserving the
competitors that set the auction's competitive outside option; set-asides buy
more reallocation by removing that outside option.

This frontier is particularly informative because the low-cost region is not a
knife-edge at 10 percent. In the preference grid, welfare costs remain inside
Monte Carlo noise up to \prefSafeBandHiNp{} percent in non-pharmaceuticals and
\prefSafeBandHiPh{} percent in pharmaceuticals, and even a 30 percent
preference costs only \welfPrefThirtyNp/\welfPrefThirtyPh{} percent of
$p_{S_1}$. The set-aside is therefore not the natural default instrument in
thick standardized procurement. It is an extreme point on the policy frontier:
maximum formal eligibility protection, purchased by deleting rival bidders from
the auction.

\subsection{How much must the planner value SME surplus?}

The distributional case for exclusion can be stated as a welfare-weight
condition. Following the social marginal welfare weight logic in
\citet{saezstantcheva2016}, let $w^{\mathrm{SME}}$ be the planner's weight on
one real of SME producer surplus relative to one real of general surplus. The
planner is indifferent between the set-aside and the 10 percent preference when
\begin{equation}
    w^{\mathrm{SME}}_\star
    =
    \frac{\mathrm{Loss}(V_0)-\mathrm{Loss}(V_3)}
         {\mathrm{SMESurplus}(V_0)-\mathrm{SMESurplus}(V_3)} .
    \label{eq:welfare_weight_v8}
\end{equation}
Because $V_3$ has near-zero welfare cost, the numerator is essentially the
deadweight cost of exclusion. The denominator is the additional SME surplus
created by moving from a price preference to a full set-aside. Under the main
specification, the implied indifference weights are
\welfWeightMainNp{} in non-pharmaceuticals and \welfWeightMainPh{} in
pharmaceuticals. In words, the planner must value an extra real of SME surplus
at more than twice a real of general surplus to prefer excluding rival bidders
over preserving them with a price preference.

The strict-invariance benchmark clarifies the scope of that statement. In
non-pharmaceuticals, the preference rule continues to dominate the set-aside;
the implied weight remains above the utilitarian benchmark. In pharmaceuticals,
the ranking is sensitive to whether the post-policy SME pool is interpreted as
a selected equilibrium participant pool or forced to have the same primitive
cost distribution as the pre-policy pool; under strict invariance the implied
weight falls to \welfWeightSiPh. That sensitivity is not a nuisance detail. It
is the policy boundary of the paper: the argument against bidder exclusion is
strongest where the protected pool is thick enough that redistribution can be
achieved without changing the price-forming bidders.

\subsection{When is exclusion defensible?}

The results do not imply that set-asides are never defensible. They imply that
the defense must be explicit. Exclusion becomes more plausible when the
protected pool is thin, when the planner's objective is to create protected
market access rather than merely tilt marginal awards, when dynamic SME
capacity gains are expected to be large, or when the planner places a high
social weight on SME surplus. Those conditions are not identified by the legal
threshold alone. They are market-structure conditions.

Non-pharmaceutical Group-65 procurement is the hard case for the set-aside: the
SME pool is thick, goods are standardized, and a price preference preserves
competition while delivering SME wins at almost no welfare cost. Pharmaceuticals
are the boundary case: the SME pool is thinner, composition changes more under
the policy, and the welfare ranking depends on how the post-policy SME
distribution is interpreted. The broader design principle is therefore
conditional rather than universal. Use exclusion only when the planner is
prepared to pay for the missing competitors; otherwise, support SMEs with an
instrument that keeps those competitors in the auction.

\section{Robustness to Core Threats}\label{sec:robustness}

The robustness checks follow the empirical chain. The legal shock must line up
with the timing of prices; drop-out bids must recover the cost objects needed
for the order-statistic comparison; the post-policy SME pool must not be the
sole source of the result; coordination must not intensify after restriction;
and the preference benchmark must not be a mechanical artifact. The question is
whether any of these threats overturns the conclusion that bidder exclusion is
the main source of the price increase.

\subsection{Reduced-form timing and scale}

The reduced-form evidence is used for timing and scale, not for the structural
decomposition. The \structuralWindowM-month benchmark in
Table~\ref{tab:did_benchmark_v8} is \didPriceEighteenPbu{} in the paper's
DiD-in-reverse convention, implying a \didLowerPct--\didUpperPct{} percent
price movement. The estimate is stable across the 6-, 12-, and 18-month
windows. Excluding the BEC enablement months before the mass-take-up cutoff
changes the coefficient from \phasedDidFull{} to \phasedDidTrunc{} with a
standard error of \phasedDidSe, and placebo cutoffs before the reform are
smaller in magnitude (\placeboSeventeenSep, \placeboSeventeenMar, and
\placeboSeventeenJun) than the real-cutoff estimates.

The conservative reading remains necessary. Group~65 is not perfectly balanced
against the pooled controls: \balanceNVarsAboveTen{} of \balanceNVarsTotal{}
pre-period covariates exceed the 0.10 normalized-difference threshold, and
\balanceNVarsAboveTwentyFive{} exceed 0.25. Modern DiD variants also reduce
precision: the BJS imputation estimate is \didAltBjsEst{} and the
Callaway--Sant'Anna estimate is \didAltCsEst. The reduced form therefore
validates the episode's sign and timing; it does not identify the exclusion and
protected-pool channels. Online Appendix~OA-B reports the balance table,
placebos, and alternative estimators.

\subsection{Cost recovery and auction primitives}

The cost-recovery concern is that the decomposition could be an artifact of
winner censoring, common auction-level shocks, or bidder dependence. The
diagnostics do not support that reading.
Alternative treatments of winner censoring give similar net price effects:
\costRegimeLosersNp/\costRegimeAllNp/\costRegimeTurnNp{} in non-pharmaceuticals
under losers-only/all-bidders/Turnbull inputs, and
\costRegimeLosersPh/\costRegimeAllPh/\costRegimeTurnPh{} in pharmaceuticals.
The Turnbull within-auction share remains large
(\turnbullShareNp{} percent in non-pharmaceuticals and \turnbullSharePh{}
percent in pharmaceuticals), so the result is not driven by treating the
winner's final bid as an exact cost observation.

The other diagnostics point the same way. The auction-level heterogeneity
correction removes common scale shocks before simulation. A Convite first-price
GPV recovery and the Preg\~ao drop-out recovery line up in the key
pharmaceutical non-SME pre-period cell after that correction. A Gaussian-copula
relaxation allows within-auction cost correlation up to $\rho_c=\ipvRhoMax$;
across that grid, the within-auction share moves by less than
\ipvShareDriftPp{} percentage points and the total effect by less than
\ipvTotalDriftPct{} percent. Tightening or loosening the
$c_\epsilon\le\filterCepsUpper$ filter and using 18-, 12-, or 6-month windows
also leave the within-auction component as the larger component in both
classes. Online Appendix~OA-C reports the cost-distribution quantiles and
primitive-stability diagnostics behind these summaries.

\subsection{Protected-pool composition}

The main modeling threat is that $S_3-S_2$ mixes additional SME bidders with a
changed post-policy SME cost distribution. The strict-invariance benchmark
sets $F_c^{\mathrm{SME,Post}}=F_c^{\mathrm{SME,Pre}}$ while keeping observed
post-policy SME entry counts. The total price effect remains positive:
\siDeltaTotalNp{} in non-pharmaceuticals and \siDeltaTotalPh{} in
pharmaceuticals. The within-auction shares rise to \strictShareNp{} and
\strictSharePh{} percent. Composition therefore matters, but it is not what
makes the exclusion component dominate the price decomposition.

Composition does matter for the policy ranking in pharmaceuticals.
Post-policy pharmaceutical SME participation is heavily composed of new firms:
\turnoverPhPctNewFirms{} percent of post-period pharma SME firms are absent
from the pre-period pool, and they account for \turnoverPhPctNewBids{} percent
of post-period pharma SME bids. In non-pharmaceuticals, new firms account for
\turnoverNpPctNewBids{} percent of post-period SME bids. The pharmaceutical
welfare ranking is therefore conditional, while the non-pharmaceutical ranking
is the more stable policy result.

\subsection{Coordination}

Coordination is the conduct threat that would most directly contaminate the
drop-out interpretation. The data show bidder clustering: Conley-style
close-pair screens \citep{conley2016} reject the independent-draw null in all
four class-by-period cells, and persistent-pair diagnostics in the spirit of
\citet{bajariye2003} flag excess repeated co-bidding. The relevant question is
whether clustering increases after non-SMEs are removed. It does not. The
Conley realized close-pair share is essentially flat in non-pharmaceuticals
(\collConleyRealPreNp{} to \collConleyRealPostNp{} percent) and falls in
pharmaceuticals (\collConleyRealPrePh{} to \collConleyRealPostPh{} percent).
The Bajari--Ye $T_1$ observed-to-null ratio also falls in both classes. Residual
baseline clustering remains a limitation, but the strongest alternative story,
a post-restriction coordination shock among surviving SMEs, is not supported by
the diagnostics. Online Appendix~OA-D reports the coordination-screen table.

\subsection{Policy benchmark}

The price-preference comparison is a static design benchmark, so the relevant
robustness checks are mainly mechanical. The non-pharmaceutical ranking remains
$V_3 \succ V_0$ over $\lambda\in[\welfLambdaRangeLo,\welfLambdaRangeHi]$ under
both the main and strict-invariance specifications; the pharmaceutical ranking
remains conditional on the treatment of the post-policy SME pool. The result is
also not an artifact of exact Vickrey equivalence. A Maskin--Riley first-price
upper-bound adjustment \citep{maskinriley2000} is at most
\maskinRileyBoundPct{} percent of the reference price, while the simulated
$V_0$--$V_3$ price gap is \vZeroVThreeGapLo--\vZeroVThreeGapHi{} percent.
Online Appendix~OA-E reports the preference grid and welfare-weight
sensitivities.

The checks leave a narrow claim. The result does not depend on the reduced-form
DiD carrying the whole design, on a single treatment of winner censoring, or on
holding the post-policy SME pool fixed. When the excluded bidders are the ones
disciplining the order statistic, set-asides are costly because they replace
competition with eligibility. The protected-pool response offsets part of that
cost, but it does not replace the competitive discipline that the policy
removes.

\section{Conclusion}\label{sec:conclusion}

This paper shifts the evaluation of SME set-asides from access alone to
price formation. A set-aside does not merely direct contracts toward
protected firms; it changes the bidder pool from which the auction price
is formed. When the excluded non-SMEs are the bidders that discipline the
relevant order statistic, SME access is purchased by weakening the
competitive mechanism itself.

The paper does not estimate a full reduced-form treatment effect of the
legal reinterpretation, nor does it solve a dynamic entry model of SME
development. It identifies a narrower object: how the price-forming order
statistic changes when an exclusionary rule removes non-SMEs and the
protected SME pool responds. Within that static auction margin, exclusion
is expensive in standardized non-pharmaceutical procurement. A planner
can still prefer set-asides, but the justification must come from
benefits outside the measured static margin: SME surplus weights,
market-access objectives, or dynamic capacity gains.

The evidence from S\~ao Paulo's procurement reform shows that the
protected pool responds, but incompletely. SME participation rises after
the expansion of SME-only tendering, yet the post-policy protected pool
does not recreate the price discipline supplied by excluded non-SMEs. In
standardized non-pharmaceutical procurement, this leaves a large static
cost: the full set-aside generates a welfare loss of
\welfLossPctLthirtyNp{} percent of the open-regime price at
$\lambda=\lambdaMain$. Pharmaceutical procurement exhibits the boundary
of the exercise. There, the protected pool is thinner, composition
changes more, and the welfare ranking becomes model-sensitive.

The policy implication is that the instrument matters. A full set-aside
is a high-redistribution instrument because it forces eligibility by
removing rival bidders. A price preference delivers less redistribution,
but it preserves the non-SME bidders that discipline the price-forming
order statistic. The relevant comparison runs not between SME support
and no support, but between exclusionary redistribution and support that
keeps the price-forming bidder pool inside the auction. The evidence
does not imply that an implemented preference regime would reproduce the
static benchmark exactly; it shows that exclusion is the costly margin,
and that preserving the price-forming pool is valuable before any
behavioral response is modeled.

In thick standardized markets, where a preference rule can tilt
allocation without deleting rival bidders, exclusion is an expensive way
to redistribute. Support for SMEs need not remove the bidders that
discipline the price.

Two extensions would sharpen the result. The drop-out decomposition can be
run on other open reverse-auction platforms---federal ComprasNet is the
natural next case---to test whether exclusion dominance travels beyond a
single state and a single product group. And linking winning SMEs to firm
panels would let the dynamic capacity gains that sit outside the static
margin here be measured directly, turning the welfare-weight threshold into
an estimate rather than an indifference point.


\section*{Declaration of competing interests}

The author declares that he has no known competing financial interests or
personal relationships that could have appeared to influence the work reported
in this paper.

\section*{Data availability}

The structural sample is derived from BEC administrative microdata accessed
through a research agreement with the Secretaria da Fazenda e Planejamento do
Estado de S\~ao Paulo. Public replication materials will be organized for the v8
version once the analysis set is frozen.

\clearpage
\bibliographystyle{chicago}
\bibliography{References}

\end{document}
